%============================================================
% BLOCK 0: subsection opening (lead-in for all V2 ablations)
% Insert as a new subsection AFTER the existing Ablation Study
% content (after the tab:rag_ablation table block, which ends
% ~line 1792) and BEFORE \subsection{Qualitative Results}
% (line 1795).
%============================================================
\subsection{Ablation Study of ReMoMask-2}
\label{subsec:ablation_v2}
The ablations above isolate the contributions of HBM, TSM, and SSTA. We now ablate the two additions of the journal version: the latent-aligned retrieval space and the re-enabled value-pathway routing.
% TODO-REF: once the methodology sections land, add \ref{} to Sec. 4.6 (latent-aligned retrieval) and Sec. 4.7 (value-pathway routing); suggested labels: subsec:latent_retrieval, subsec:rt_value.
Unless stated otherwise, all variants share the identical backbone, masking strategy, and training configuration, and are evaluated on HumanML3D with the complete masked--residual generation pipeline under the same protocol as Table~\ref{tab:t2m_experiment}.

%============================================================
% BLOCK 1: table-orthogonal (MAIN 2x2 ablation)
%============================================================
\begin{table}[t]
    \centering
    \caption{\textbf{Orthogonal ablation of retrieval space and value-pathway routing on HumanML3D.} All four variants share the identical backbone, masking strategy, and training configuration; only the retrieval space (semantic HBM space vs.\ latent-aligned $z_e$ space) and the routing of the retrieved text embedding $R_t$ into the Value pathway vary. The first row corresponds to the conference ReMoMask configuration retrained under the journal protocol; the last row is the full ReMoMask-2.}
    \label{tab:ablation_v2_orthogonal}
    \renewcommand{\arraystretch}{1.3}  % Adjust the line height (global)
    \resizebox{\columnwidth}{!}{
        \begin{tabular}{l| c| c| c| c}
            \toprule
            Retrieval Space & $R_t$ in Value & FID$\downarrow$ & Top1$\uparrow$ & MMDIST$\downarrow$ \\
            \hline
            \hline
            Semantic (HBM) & $\times$ & $x.xxx^{\pm.xxx}$ & $x.xxx^{\pm.xxx}$ & $x.xxx^{\pm.xxx}$ \\
            Semantic (HBM) & $\checkmark$ & $x.xxx^{\pm.xxx}$ & $x.xxx^{\pm.xxx}$ & $x.xxx^{\pm.xxx}$ \\
            Latent-aligned ($z_e$) & $\times$ & $x.xxx^{\pm.xxx}$ & $x.xxx^{\pm.xxx}$ & $x.xxx^{\pm.xxx}$ \\
            \rowcolor{yellow!20} Latent-aligned ($z_e$) & $\checkmark$ & $x.xxx^{\pm.xxx}$ & $x.xxx^{\pm.xxx}$ & $x.xxx^{\pm.xxx}$ \\
            \bottomrule
        \end{tabular}
    }
\end{table}

%============================================================
% BLOCK 2: table-alignment-objective (KL vs InfoNCE)
%============================================================
\begin{table}[t]
    \centering
    \caption{Alignment objective for the query projector $\phi$ on HumanML3D. Retrieval columns report text-to-motion recall of the projected query in the $z_e$ space; FID reports downstream generation quality when the resulting projector is used in ReMoMask-2. The adopted setting is shaded.}
    \label{tab:ablation_v2_alignment}
    \renewcommand{\arraystretch}{1.3}  % Adjust the line height (global)
    \resizebox{\columnwidth}{!}{
        \begin{tabular}{l| ccc| c}
            \toprule
            Objective & R@1$\uparrow$ & R@5$\uparrow$ & R@10$\uparrow$ & FID$\downarrow$ \\
            \hline
            \hline
            InfoNCE & xx.xx & xx.xx & xx.xx & $x.xxx^{\pm.xxx}$ \\
            \rowcolor[gray]{0.90} KL distillation & xx.xx & xx.xx & xx.xx & $x.xxx^{\pm.xxx}$ \\
            \bottomrule
        \end{tabular}
    }
\end{table}

%============================================================
% BLOCK 3: table-teacher (HBM vs TMR teacher)
%============================================================
\begin{table}[t]
    \centering
    \caption{Choice of the distillation teacher for the query projector $\phi$ on HumanML3D. The projector is distilled from each teacher's top-256 soft targets under otherwise identical settings. The adopted setting is shaded.}
    \label{tab:ablation_v2_teacher}
    \renewcommand{\arraystretch}{1.3}  % Adjust the line height (global)
    \resizebox{\columnwidth}{!}{
        \begin{tabular}{l| ccc| c}
            \toprule
            Teacher & R@1$\uparrow$ & R@5$\uparrow$ & R@10$\uparrow$ & FID$\downarrow$ \\
            \hline
            \hline
            TMR~\cite{tmr} & xx.xx & xx.xx & xx.xx & $x.xxx^{\pm.xxx}$ \\
            \rowcolor[gray]{0.90} HBM & xx.xx & xx.xx & xx.xx & $x.xxx^{\pm.xxx}$ \\
            \bottomrule
        \end{tabular}
    }
\end{table}

%============================================================
% BLOCK 4: table-sensitivity (distillation Top-K + projector capacity)
% Retrieval-level only (no FID column) to keep the sweep cheap and compact.
%============================================================
\begin{table}[t]
    \centering
    \caption{Sensitivity of the query projector to the number of distillation targets (Top-$K$; $K$ here denotes the number of teacher candidates, not body parts) and to projector capacity, measured by text-to-motion recall on HumanML3D. Default settings are shaded.}
    \label{tab:ablation_v2_sensitivity}
    \renewcommand{\arraystretch}{1.3}  % Adjust the line height (global)
    \resizebox{\columnwidth}{!}{
        \begin{tabular}{l| c| ccc}
            \toprule
             & \#Params & R@1$\uparrow$ & R@5$\uparrow$ & R@10$\uparrow$ \\
            \hline
            \hline
            \multicolumn{5}{l}{\emph{Distillation Top-$K$}} \\
            \hline
            $K{=}64$ & 1.57M & xx.xx & xx.xx & xx.xx \\
            \rowcolor[gray]{0.90} $K{=}256$ (default) & 1.57M & xx.xx & xx.xx & xx.xx \\
            $K{=}1024$ & 1.57M & xx.xx & xx.xx & xx.xx \\
            \hline
            \multicolumn{5}{l}{\emph{Projector capacity}} \\
            \hline
            Linear ($512{\to}1024$) & 0.53M & xx.xx & xx.xx & xx.xx \\
            \rowcolor[gray]{0.90} MLP, hidden 1024 (default) & 1.57M & xx.xx & xx.xx & xx.xx \\
            MLP, hidden 2048 & 3.15M & xx.xx & xx.xx & xx.xx \\
            \bottomrule
        \end{tabular}
    }
\end{table}

%============================================================
% BLOCK 5: prose-analysis (dummy analysis, written as-if-expected;
% every placeholder-dependent claim marked \textbf{[TBD]})
%============================================================
\noindent \myparagraph{Retrieval Space and Value Routing.}
Table~\ref{tab:ablation_v2_orthogonal} disentangles the two journal-version factors. The retrieval space is the dominant one: switching from the semantic HBM space to the latent-aligned $z_e$ space reduces FID by \textbf{[TBD]} without routing and by \textbf{[TBD]} with routing, together with consistent gains in Top-1 R-Precision. Routing $R_t$ into the Value pathway contributes a smaller gain in both retrieval spaces (\textbf{[TBD]} and \textbf{[TBD]} FID reduction, respectively), and the two effects compose nearly additively: the full configuration ($z_e$ + routing) attains the best results across all three metrics, indicating that the two contributions are largely independent.

The routing result should be read against the premise it rests on. In the conference version, $R_t$ was a text-domain signal from the contrastive semantic space, and injecting it into the Value pathway degraded quality (FID $0.104$ vs.\ $0.027$ in Table~\ref{tab:ablation}), which motivated excluding textual semantics from Value to prevent domain mismatch. In ReMoMask-2, $R_t$ is projected by $\phi$ into the motion-aligned $z_e$ space and thus enters SSTA as a motion-domain signal; the premise behind the exclusion no longer applies to it, and re-enabling the routing is accordingly beneficial, with the clearly larger gain observed in the latent-aligned rows (\textbf{[TBD]}). The modest positive effect in the semantic-space rows (\textbf{[TBD]}) suggests that the routing penalty observed at the conference operating point does not carry over unchanged to the journal re-evaluation; it does not overturn the conference finding, which isolates a semantic-space $R_t$ under the original protocol.

\noindent \myparagraph{Alignment Objective.}
Table~\ref{tab:ablation_v2_alignment} compares KL distillation with InfoNCE for training $\phi$. KL distillation yields both higher projector recall and lower downstream FID (\textbf{[TBD]}). This is consistent with the distillation view of the alignment problem: the projector does not need to learn a text--motion metric from scratch, but to reproduce, inside the $z_e$ space, the ranking of an existing stronger retriever. The top-256 soft targets convey the teacher's full similarity structure over candidates, whereas InfoNCE reduces supervision to a single positive per caption; since the database expands 23,384 motions into 66,912 (motion, caption) entries, roughly three captions per motion, a single-positive objective can treat entries of the very same motion as negatives, injecting label noise that the soft-target objective avoids.
% TODO-CITE: Hinton et al., "Distilling the Knowledge in a Neural Network" (arXiv:1503.02531) — canonical reference for KL soft-target distillation.
% TODO-CITE: van den Oord et al., "Representation Learning with Contrastive Predictive Coding" (arXiv:1807.03748) — original InfoNCE objective.

\noindent \myparagraph{Distillation Teacher.}
Table~\ref{tab:ablation_v2_teacher} compares distilling $\phi$ from the HBM retriever against distilling from TMR~\cite{tmr}. The HBM teacher yields better projector recall and better downstream generation (\textbf{[TBD]}). This mirrors the retrieval gap between the teachers themselves (18.49 vs.\ 5.68 text-to-motion R@1 on HumanML3D in Table~\ref{tab:rag_experiment}): since the projector learns the teacher's ranking, its quality is effectively upper-bounded by the teacher's ranking quality, and improvements to the retriever transfer directly to the latent-aligned pipeline. This also links the two versions of our system: the conference contribution (HBM) is what makes the journal contribution (latent-aligned retrieval) train well.

\noindent \myparagraph{Geometry of the $z_e$ Space.}
The measured geometry of the pooled $z_e$ embeddings explains both why latent retrieval works and why a learned projector is required. Pairwise cosine similarities average $0.62$ with a standard deviation of $0.25$ (range $[-0.55, 0.998]$)---far from collapsed---so cosine similarity remains discriminative and nearest-neighbor retrieval in $z_e$ is well-posed. At the same time, the space is severely anisotropic: only about $11$ of $1024$ dimensions are effective, and $7$ principal directions explain $95\%$ of the variance. Usable structure therefore occupies a narrow low-dimensional cone of the ambient space; a fixed or analytically constructed mapping (e.g., reusing the semantic-space text embedding directly) would spread the query's energy across mostly uninformative coordinates and destroy the ranking, whereas a projector distilled from teacher rankings learns to place queries inside precisely this cone. The same anisotropy favors ranking-level (KL) supervision over pointwise regression to motion embeddings, consistent with Table~\ref{tab:ablation_v2_alignment}.
% TODO-REF: add Fig.~\ref{fig:ze_geometry} (cosine histogram + eigenvalue spectrum) reference here once the Design Principles addition introducing fig:ze_geometry is merged; keep this sentence out until the label exists.

\noindent \myparagraph{Distillation Top-$K$ and Projector Capacity.}
Table~\ref{tab:ablation_v2_sensitivity} probes the two main hyper-parameters of the alignment stage. Performance is stable in a broad range around the default $K{=}256$: a very small candidate set truncates the teacher's ranking signal, while a very large one dilutes it with a low-similarity tail (\textbf{[TBD]}). Increasing projector capacity beyond the default two-layer MLP brings no consistent gain (\textbf{[TBD]}), while a single linear map underfits the teacher ranking (\textbf{[TBD]}). Both trends indicate that the improvement of ReMoMask-2 stems from \emph{where} the queries are aligned---the generator's own latent space---rather than from added capacity in the alignment module.

%============================================================
% BLOCK 6: implementation-details-addition
% Append at the END of \subsection{Implementation Details},
% right after the sentence "All models are implemented in PyTorch."
% (line 1548). Does not modify or contradict the existing
% "8 Tesla A800" sentence: the V2 paragraph fixes the *configuration*,
% not the hardware, for the controlled comparison.
%============================================================
\noindent \myparagraph{ReMoMask-2.}
The $z_e$ retrieval database is constructed once from the frozen 2D-RVQ-VAE: each of the 23,384 training motions is passed through the 2D encoder, and its pre-quantization latent (1024 channels over the temporal--part grid) is mean-pooled over the temporal and part axes and $\ell_2$-normalized to form a 1024-dimensional retrieval key; expanding motions to their captions yields 66,912 (motion, caption) entries, and the full build takes about five minutes. The query projector $\phi$ is a two-layer MLP (Linear $512{\to}1024$, GELU, Linear $1024{\to}1024$; $\approx$1.57M parameters) applied to the same frozen CLIP~\cite{clip} ViT-B/32 text embedding used by the retrieval model, with $\ell_2$-normalized outputs. It is trained for 200 epochs with a batch size of 128 on a single GPU using Adam with cosine annealing, distilling the HBM retriever's top-256 soft targets at temperature $\tau = 0.07$; the VQ-VAE remains frozen throughout. On the generator side, the SSTA architecture is unchanged: the only additional parameters are two linear adapters ($1024 \to 512$) that project the retrieved $R_m$ and $R_t$ to the model dimension. For the controlled comparison in this journal version, the masked models of ReMoMask and ReMoMask-2 are trained under exactly the same configuration---identical architecture, optimization hyper-parameters, batch size, training schedule, and random seed---so that the retrieval module is the only varying factor between the two systems.